\documentclass[12pt,a4paper,twoside]{article}
\usepackage{tabularx}
\usepackage[utf8]{inputenc}
\usepackage[galician]{babel}
\usepackage{amsmath}
\usepackage[usenames,dvipsnames,svgnames,table]{xcolor}
\usepackage{amsfonts}
\usepackage{amssymb}
\usepackage{upgreek}
\usepackage{booktabs}
\usepackage{esint}
\usepackage{fancybox}
\usepackage{graphicx}
\usepackage{multirow}
\usepackage{subcaption}
\usepackage{float}
\usepackage[letterpaper,top=3.5cm,bottom=2cm,left=3cm,right=3cm,marginparwidth=1.75cm,footskip=1.25cm %, showframe
]{geometry}
\usepackage[hidelinks]{hyperref}
\usepackage{pgfplots}
\usepackage{multicol}
\usepackage{braket}
\usepackage{upgreek}
\usepackage{tensor}
\usepackage{changepage}

\usepackage{diffcoeff}[=v4]
\diffdef{ pvrule } { op-symbol = \partial}

\usepackage{titlesec}
% Teño que ver mellor como facer con este paquete para crear unha clase menor a \subsubsection{}
%\titleformat{\part}[hang]{\normalfont\Huge\bfseries}{\texttt{\thepart} -}{0.5em}{}
\titleclass{\subsubsub}{straight}[\subsubsection] 
\newcounter{subsubsub} 
\renewcommand{\thesubsubsub}{\Alph{subsubsub}}

%\titlecontents{\subsubsub}[⟨left⟩]{⟨above-code⟩} {⟨numbered-entry-format⟩}{⟨numberless-entry-format⟩} {⟨filler-page-format⟩}[⟨below-code⟩]

% Pódeselle meter unha tipografía guapa
\author{Sebastián Táboas Pazo\\
35679654F}
\title{\huge{\textbf{Teorías de Lovelock}}}
\date{}

% Fáltame implementar a data na marxe superior (por defecto xenérase inmediatamente despois da autoría)

\usepackage{fancyhdr,tikz}
\renewcommand{\headrulewidth}{1.5pt}


\makeatletter
\renewcommand{\footrule}{
    \if@fancyplain\let\footrulewidth\plainfootrulewidth\fi
      \ifdim\footrulewidth>0pt
        \vspace{20pt}
        \makebox[\textwidth]{
          \rule[1ex]{0.46\textwidth}{1.6pt}
          \hspace{1em}
          \raisebox{0.5ex}{\thepage}
          \hspace{1em}
          \rule[1ex]{0.46\textwidth}{1.6pt}
        }
        \vspace{6pt}
      \fi
    }
\makeatother

\renewcommand{\sectionmark}[1]{
  \markboth{\thesection. - \MakeUppercase{#1}}{}
}

\NewDocumentCommand\Cita{ m m }{%
    \begin{minipage}[t]{\hsize}
        \vspace*{0.5em}
        \begin{quotation}
            \begin{flushright}
            {\noindent \itshape #1}\\
            \end{flushright}
            \vspace{-4ex}
            % \hfill — #2 % Como estaba antes <-
            \hfill
            \begin{flushright}
            — #2
            \end{flushright}
        \end{quotation}
        \vspace*{0.5em}
    \end{minipage}
}

\fancypagestyle{main}{
    \fancyhf{}
    \fancyhead[RO]{\bfseries\leftmark}
    \fancyhead[LE]{\texttt{\rightmark}}
    \fancyhead[RE]{Sebastián Táboas Pazo}
    \renewcommand{\footrulewidth}{1.5pt}
    %\fancyfoot[C]{\hskip-0.01cm\vtop{\vskip.05cm\hsize=7.5cm \drulefill}
    % ~~\thepage ~~\hskip0.01cm \vtop{\vskip.05cm\hsize=7.44cm \drulefill}}
}

\fancypagestyle{anexo}{
    \fancyhf{}
    \fancyhead[RO]{\textbf{Anexos -}}
    \fancyhead[LE]{\texttt{\rightmark}}
    \renewcommand{\footrulewidth}{1.5pt}
    %\fancyfoot[C]{\hskip-0.01cm\vtop{\vskip.05cm\hsize=7.5cm \drulefill}
    % ~~\thepage ~~\hskip0.01cm \vtop{\vskip.05cm\hsize=7.44cm \drulefill}}
}


\fancypagestyle{index}{
    \fancyhf{}
    \rhead[Gravitación - 2025/2026]{}
    \lhead[]{\rightmark}
    \renewcommand{\footrulewidth}{0pt}
}

\fancypagestyle{bib}{
    \fancyhf{}
    \fancyhead[LE]{\bfseries\leftmark}
    \fancyhead[RO]{Gravitación - 2025/2026}
    \fancyhead[RE]{}
    \renewcommand{\footrulewidth}{1.5pt}
}
\fancypagestyle{apaño1}{
    \fancyhf{}
    \fancyhead[LE]{\texttt{5.1.\hspace{1em}Teorema de Lovelock}}
    \fancyhead[RE]{Sebastián Táboas Pazo}
    \renewcommand{\footrulewidth}{1.5pt}
    %\fancyfoot[C]{\hskip-0.01cm\vtop{\vskip.05cm\hsize=7.5cm \drulefill}
    % ~~\thepage ~~\hskip0.01cm \vtop{\vskip.05cm\hsize=7.44cm \drulefill}}
}\fancypagestyle{apaño2}{
    \fancyhf{}
    \fancyhead[LE]{\texttt{5.2.\hspace{1em}Acción de Lovelock en 2, 4 e 5 dimensións}}
    \fancyhead[RE]{Sebastián Táboas Pazo}
    \renewcommand{\footrulewidth}{1.5pt}
    %\fancyfoot[C]{\hskip-0.01cm\vtop{\vskip.05cm\hsize=7.5cm \drulefill}
    % ~~\thepage ~~\hskip0.01cm \vtop{\vskip.05cm\hsize=7.44cm \drulefill}}
}
\usepackage{titling}
\setlength{\droptitle}{2.5cm}

% estos son os comandos que poñen taboaa, figura, ec como (s.n)
% 
% \makeatletter
% \@addtoreset{equation}{section}
% \renewcommand{\theequation}{\arabic{section}.\arabic{equation}}
% \def\equation@reset{%
%   \setcounter{equation}{-1}} % Comezar dende -1 para que a primeira sexa 0
% \makeatother

% \makeatletter
% \@addtoreset{table}{section}
% \renewcommand{\thetable}{\arabic{section}.\arabic{table}}
% \def\table@reset{%
%   \setcounter{table}{-1}} % Comezar dende -1 para que a primeira sexa 0
% \makeatother

% \makeatletter
% \@addtoreset{figure}{section}
% \renewcommand{\thefigure}{\arabic{section}.\arabic{figure}}
% \def\figure@reset{%
%   \setcounter{figure}{-1}} % Comezar dende -1 para que a primeira sexa 0
% \makeatother

\usepackage[backend=biber,style=alphabetic,sorting=nyt]{biblatex}

\addbibresource{./bibliografia.bib}
 

\begin{document}

\pagenumbering{gobble}
\begin{titlepage}
    \centering
    \maketitle
    \vspace{20pt}
    \begin{tikzpicture}[baseline, every node/.style={anchor=center}]
        \begin{axis}[hide axis,scale=0.4,anchor=center,axis equal,
                width=33cm,height=12cm]
            \addplot[black, thick, smooth] table
            %[x expr=\thisrowno{0}*1.5, y expr=\thisrowno{1}*1.5]
                {./clotoide.dat}; % Sin puntos y con curva suave
            \addplot[thick, black, domain=-3:-1, samples=100] plot (\x, {0});
            \addplot[thick, black, domain=1:3, samples=100] plot (\x, {0});
        \end{axis}
    \end{tikzpicture}
    \vfill
    \begin{minipage}[t]{1\textwidth}
    \centering
    \includegraphics[width=0.56\textwidth]{normal.png}
    \end{minipage}
    \vspace{1.6cm}\newline
    \texttt{Data: \today}
\end{titlepage}\newpage%\null\newpage

\pagestyle{index}
\tableofcontents
\newpage

\pagestyle{main}
\pagenumbering{arabic}

\section*{Índice de contidos}
\begin{itemize}
	\item 0. Introdución
	\item 1. Depuración de datos
	\item 2. Representacións estatísticas
	\item 3. Análise con sklearn (regresión loxística)
	\item \hspace{4em}3.0.1 Modelo Linear
	\item \hspace{4	em}3.0.2 Modelo RBF
	\item \hspace{2em}3.1 Análise sen sesgo
	\item \hspace{4em}3.1.1 Estimación, training e resultados do Modelo Linear/Loxístico linear
	\item \hspace{4em}3.1.2 Estimación, training e resultados do Modelo RBF
	\item \hspace{4em}3.1.3 Estimación, training e resultados do Modelo RandomForest
	\item \hspace{4em}3.1.4 Estimación, training e resultados do Modelo HistGradientBoosting
	\item \hspace{2em}3.2 Conclusións
	\item \hspace{2em}3.3 Análise de regresións (Score)
	\item 4. Análise con redes neuronais
	\item 5. Conclusións finais
	\item Apéndice Cross Validation en redes neuronais
\end{itemize}

\section*{Resumo}
O noso obxectivo neste traballo foi predicir resultados dun dataframe mediante métodos de machine learning. Primeiro limpiamos un pouco os nosos datos e facemos e os adaptamos para podelos procesor. Nun segundo paso estudamos un pouco como son esos datos aínda que non nos paramos realmente a facer ningunha clase de estudo estatístico. Despois empregamos funcións da librería de sklearn para obter predicións. Nun primer momento aplicamos o modelo máis sinxelo, unha regresión linear que optimizamos os seus parámetros cunha cross validation pero máis desglosada e mostrando as vísceras do cálculo; facelo así supón unha maior carga computacional pero vese mellor o que fan os programas de cross validation que usaremos máis adiante con GridSearchCV ou keras_tuner.RandomSearch. Os resultados da recta axustada son aínda bastante pésimos, empregamos tamén un modelo RBF (Radial Basis Function) que considera no linearidades mellorando increíblemente o sistema. Como atopamos que este segundo ten unha predición altísima consideramos que hai algunha das nosas variables que está moi correlacionada coa que queremos determinar aplicando mutual_info_classif.

Volvemos a crear modelos pero esta vez antes de aplicalos realizamos validación cruzada coa función GridSearchCV para optimizar os parámetros dos modelos. Nun primeiro lugar repetimos co modelo linear (loss = 'hinge') que xa vemos que non nos dá moi bos resultados e o comparamos coa regresión loxística (loss = 'log_loss) e xa por aforrar só vemos que tal procesa a regresión loxística. Aplicamos cv sobre a rbf tamén e seguimos obtendo moi bos resultados (de feito será a técnica que mellor prediga os nosos resultados durante todo o traballo). Engadimos agora dous modelos máis sobre os que tamén facemos cv: RandomForest e HistGradientBoosting. O primeiro dá resultados mellores que a regresión loxística pero cae 5 puntos detrás do rbf e ademais ten overfitting (isto vese en que a accuracy_score para os datos de entrenamento ronda o valor da unidade). O segundo modelo mellora respecto do anterior e, aínda que menos, segue tendo bastante overfitting.

falar das redes

\section*{Algúns resultados da optimización de redes}

Algúns resultados teñen un código identificativo
% poñer as imaxes e facer o itematize dos indices

\clearpage
\pagestyle{bib}
\renewcommand{\thesection}{}
\printbibliography[heading=bibintoc,title={Bibliografía}]

\end{document}

% Exemplos elementais do paquete 'diffcoeff'
\begin{gather}
    \diff[n]{f}{x}[x_0]\\
    \diff*{(ax^2+bx+c)}x\\
    \diffp!{F(x,y)}x\\
    \diffp ST[V]\\
    \diffp L{q_k}-\diff*{\diffp L{\dot{q}_k}[]}t = 0\\
    \diffp L{q_k}-\diff*{\diffp L{\dot{q}_k}[]}t = 0\\
    \diffp[2,3]F{x,y,z}\\
    \diff.pvrule.{F(x,y)}x[x=1]
\end{gather}
